\chapter{GPU Sorting}\label{ch:gpu-sorting}

\section{Why sorting is the linchpin}

Every rasterizing pathway composites translucent tets with premultiplied
alpha, which is only correct in depth order. rmeshvk's tets come from a
Delaunay tetrahedralization, and \code{RADIANCE_MESHES.md} records the
property everything rests on: the back-to-front order from an origin $o$ is
the descending order of the \emph{circumsphere power}
\[
  P(T) = \lVert C(T) - o \rVert^2 - r(T)^2 ,
\]
independent of view direction (Karasick et al.\ 1997; cf.\ Edelsbrunner
1989, Max 1990). One sort per origin therefore serves any projection from
that origin --- the six faces of a shadow cubemap being the notable
\emph{non}-example, since each face pass still needs its own
\emph{culling}, and compositing correctness per face requires sorting with
that face's projection (see Chapter~\ref{ch:shadow-rendering}).

\paragraph{Key construction.} Expanding $P$ around a tet vertex $v_0$ gives
the cancellation-stable form used in \code{project_compute.wgsl}:
\[
  \text{depth} = (o - v_0) \cdot (o + v_0 - 2C) ,
\]
clamped to $\pm 10^{20}$, bitcast to \code{u32}, and bit-inverted so an
\emph{ascending} radix sort yields \emph{back-to-front}. Culled slots get
\code{0xFFFFFFFF} and sort to the end; the value array carries tet ids.
Degenerate (near-planar) tets get $r^2 = 10^{20}$ at load time, sorting them
last. A 16-bit linear quantization of the key
(\code{project_compute_16bit.wgsl}, \code{SortMode::Gpu16}) halves the radix
passes when precision allows.

\section{Backend selection}

\code{rmesh-sort} picks a backend at startup:
\begin{itemize}
\item \textbf{DRS} when \code{Features::SUBGROUP} is available (native
  Vulkan/Metal/DX12);
\item \textbf{Basic} otherwise (portable, workgroup atomics only);
\item \textbf{CPU} for \code{SortMode::Cpu} and the web viewer, where wgpu
  28's WebGPU backend exposes no subgroups and the Basic sort is visibly
  out of order at multi-million tet counts.
\end{itemize}
\code{record_radix_sort} returns \code{result_in_b}: the sort ping-pongs
between A/B buffer pairs, and an odd pass count leaves the result in B, so
downstream passes carry two prebuilt bind groups and pick by this flag.

\section{Basic backend: 5-pass, 4-bit radix}

\code{crates/rmesh-sort/src/basic.rs}, with one shader per pass
(\code{radix_sort_count.wgsl} through \code{radix_sort_scatter.wgsl});
256-thread workgroups, 4 elements/thread (blocks of 1024), 16 bins. Per 4-bit digit
(8 passes for 32-bit keys):
\begin{enumerate}
\item \textbf{count} --- per-block digit histograms via workgroup atomics;
\item \textbf{reduce} --- tree-reduce each bin's block counts into a spine;
\item \textbf{scan} --- single-workgroup exclusive prefix sum of the spine;
\item \textbf{scan\_add} --- add the scanned spine back into the per-block
  counts, yielding global bin offsets;
\item \textbf{scatter} --- a stable in-block 2-bit-at-a-time local sort,
  then scatter keys and values to their global positions.
\end{enumerate}
A \code{KEY_STRIDE} substitution reuses the same shaders for 32-bit tet keys
and 64-bit tile keys.

\section{DRS backend: 3-kernel, 8-bit radix}

\code{crates/rmesh-sort/src/drs.rs} ports b0nes164's
\emph{DeviceRadixSort} (shaders \code{drs_upsweep.wgsl},
\code{drs_scan.wgsl}, \code{drs_downsweep.wgsl}): radix 256, 15
keys/thread, 3840-element partitions, 4 passes for 32 bits.
Upsweep builds per-block and global digit histograms; a 256-workgroup scan
(one per digit) prefixes the spine; the downsweep ranks keys cooperatively
with subgroup operations and scatters. The per-pass key count is copied from
a GPU counter, so the visible-tet count never touches the CPU.

\section{CPU fallback}

\code{crates/rmesh-sort/src/cpu.rs} (\code{CpuSorter}) performs
Gribb--Hartmann sphere-vs-frustum culling (wgpu's $[0,1]$ NDC makes the near
plane row $r_2$, not $r_3 + r_2$), computes the same power key with the
bit-inversion trick, packs \code{(key << 32) | tet_id} into \code{u64}, and
uses \code{sort_unstable} (pdqsort --- in place, avoiding a
$\sim$38\,MB radix scratch). The viewer uploads the resulting compact index
list and visible count with a \code{StagingBelt} between the projection and
interval-generation submits.

\section{Tile-pair sorting for the trainable path}

Training needs each screen tile's tets in depth order, which is a sort of
\emph{(tile, tet)} pairs under a 64-bit key (tile id in the high half, depth
in the low half). \code{rmesh-tile} builds the pairs without CPU readback
(\code{record_scan_tile_pipeline}):

\begin{enumerate}
\item \textbf{tile\_fill} --- fill the key/value arrays with the
  \code{0xFFFFFFFF} sentinel;
\item \code{project_compute} counts, per visible tet, the tiles its
  silhouette touches (conservative scanline test) into
  \code{tiles_touched};
\item \textbf{prepare\_dispatch} --- one thread converts the visible count
  into dispatch arguments for the following passes;
\item \textbf{RTS prefix scan} (\code{rts.wgsl}, a port of b0nes164's
  reduce-then-scan, subgroup-based, 256 threads $\times$ vec4$\times$4
  elements) turns the per-tet tile counts into write offsets --- no atomics;
\item \textbf{tile\_gen\_scan} --- each tet re-runs the scanline fill and
  emits one (key, tet id) pair per covered tile; the last tet writes the
  total pair count;
\item \textbf{radix sort} the pairs (\code{KEY_STRIDE=2}; the bit count
  adapts to the number of tile bits);
\item \textbf{tile\_ranges} --- one thread per pair detects tile boundaries
  in the sorted array and writes \code{tile_ranges[tile] = [start, end)}.
\end{enumerate}

The tiled rasterizer and backward kernels then walk \code{tile_ranges}
front-to-back and back-to-front respectively (Chapter~\ref{ch:tile-renderer}).
The forward interval renderer sets \code{tile_size_u = 0} to skip the
tile-touch counting entirely --- it is the dominant cost of the full
projection kernel and only training needs it.

